% Options for packages loaded elsewhere
\PassOptionsToPackage{unicode}{hyperref}
\PassOptionsToPackage{hyphens}{url}
\PassOptionsToPackage{dvipsnames,svgnames,x11names}{xcolor}
%
\documentclass[
  letterpaper,
]{tufte-handout}
\usepackage{amsmath,amssymb}
\usepackage{iftex}
\ifPDFTeX
  \usepackage[T1]{fontenc}
  \usepackage[utf8]{inputenc}
  \usepackage{textcomp} % provide euro and other symbols
\else % if luatex or xetex
  \usepackage{unicode-math} % this also loads fontspec
  \defaultfontfeatures{Scale=MatchLowercase}
  \defaultfontfeatures[\rmfamily]{Ligatures=TeX,Scale=1}
\fi
\usepackage{lmodern}
\ifPDFTeX\else
  % xetex/luatex font selection
\fi
% Use upquote if available, for straight quotes in verbatim environments
\IfFileExists{upquote.sty}{\usepackage{upquote}}{}
\IfFileExists{microtype.sty}{% use microtype if available
  \usepackage[]{microtype}
  \UseMicrotypeSet[protrusion]{basicmath} % disable protrusion for tt fonts
}{}
\makeatletter
\@ifundefined{KOMAClassName}{% if non-KOMA class
  \IfFileExists{parskip.sty}{%
    \usepackage{parskip}
  }{% else
    \setlength{\parindent}{0pt}
    \setlength{\parskip}{6pt plus 2pt minus 1pt}}
}{% if KOMA class
  \KOMAoptions{parskip=half}}
\makeatother
\usepackage{xcolor}
\usepackage{listings}
\newcommand{\passthrough}[1]{#1}
\lstset{defaultdialect=[5.3]Lua}
\lstset{defaultdialect=[x86masm]Assembler}
\usepackage{color}
\usepackage{fancyvrb}
\newcommand{\VerbBar}{|}
\newcommand{\VERB}{\Verb[commandchars=\\\{\}]}
\DefineVerbatimEnvironment{Highlighting}{Verbatim}{commandchars=\\\{\}}
% Add ',fontsize=\small' for more characters per line
\usepackage{framed}
\definecolor{shadecolor}{RGB}{255,255,255}
\newenvironment{Shaded}{\begin{snugshade}}{\end{snugshade}}
\newcommand{\AlertTok}[1]{\textcolor[rgb]{0.75,0.01,0.01}{\textbf{\colorbox[rgb]{0.97,0.90,0.90}{#1}}}}
\newcommand{\AnnotationTok}[1]{\textcolor[rgb]{0.79,0.38,0.79}{#1}}
\newcommand{\AttributeTok}[1]{\textcolor[rgb]{0.00,0.34,0.68}{#1}}
\newcommand{\BaseNTok}[1]{\textcolor[rgb]{0.69,0.50,0.00}{#1}}
\newcommand{\BuiltInTok}[1]{\textcolor[rgb]{0.39,0.29,0.61}{\textbf{#1}}}
\newcommand{\CharTok}[1]{\textcolor[rgb]{0.57,0.30,0.62}{#1}}
\newcommand{\CommentTok}[1]{\textcolor[rgb]{0.54,0.53,0.53}{#1}}
\newcommand{\CommentVarTok}[1]{\textcolor[rgb]{0.00,0.58,1.00}{#1}}
\newcommand{\ConstantTok}[1]{\textcolor[rgb]{0.67,0.33,0.00}{#1}}
\newcommand{\ControlFlowTok}[1]{\textcolor[rgb]{0.12,0.11,0.11}{\textbf{#1}}}
\newcommand{\DataTypeTok}[1]{\textcolor[rgb]{0.00,0.34,0.68}{#1}}
\newcommand{\DecValTok}[1]{\textcolor[rgb]{0.69,0.50,0.00}{#1}}
\newcommand{\DocumentationTok}[1]{\textcolor[rgb]{0.38,0.47,0.50}{#1}}
\newcommand{\ErrorTok}[1]{\textcolor[rgb]{0.75,0.01,0.01}{\underline{#1}}}
\newcommand{\ExtensionTok}[1]{\textcolor[rgb]{0.00,0.58,1.00}{\textbf{#1}}}
\newcommand{\FloatTok}[1]{\textcolor[rgb]{0.69,0.50,0.00}{#1}}
\newcommand{\FunctionTok}[1]{\textcolor[rgb]{0.39,0.29,0.61}{#1}}
\newcommand{\ImportTok}[1]{\textcolor[rgb]{1.00,0.33,0.00}{#1}}
\newcommand{\InformationTok}[1]{\textcolor[rgb]{0.69,0.50,0.00}{#1}}
\newcommand{\KeywordTok}[1]{\textcolor[rgb]{0.12,0.11,0.11}{\textbf{#1}}}
\newcommand{\NormalTok}[1]{\textcolor[rgb]{0.12,0.11,0.11}{#1}}
\newcommand{\OperatorTok}[1]{\textcolor[rgb]{0.12,0.11,0.11}{#1}}
\newcommand{\OtherTok}[1]{\textcolor[rgb]{0.00,0.43,0.16}{#1}}
\newcommand{\PreprocessorTok}[1]{\textcolor[rgb]{0.00,0.43,0.16}{#1}}
\newcommand{\RegionMarkerTok}[1]{\textcolor[rgb]{0.00,0.34,0.68}{\colorbox[rgb]{0.88,0.91,0.97}{#1}}}
\newcommand{\SpecialCharTok}[1]{\textcolor[rgb]{0.24,0.68,0.91}{#1}}
\newcommand{\SpecialStringTok}[1]{\textcolor[rgb]{1.00,0.33,0.00}{#1}}
\newcommand{\StringTok}[1]{\textcolor[rgb]{0.75,0.01,0.01}{#1}}
\newcommand{\VariableTok}[1]{\textcolor[rgb]{0.00,0.34,0.68}{#1}}
\newcommand{\VerbatimStringTok}[1]{\textcolor[rgb]{0.75,0.01,0.01}{#1}}
\newcommand{\WarningTok}[1]{\textcolor[rgb]{0.75,0.01,0.01}{#1}}
\setlength{\emergencystretch}{3em} % prevent overfull lines
\providecommand{\tightlist}{%
  \setlength{\itemsep}{0pt}\setlength{\parskip}{0pt}}
\setcounter{secnumdepth}{-\maxdimen} % remove section numbering
% definitions for citeproc citations
\NewDocumentCommand\citeproctext{}{}
\NewDocumentCommand\citeproc{mm}{%
  \begingroup\def\citeproctext{#2}\cite{#1}\endgroup}
\makeatletter
 % allow citations to break across lines
 \let\@cite@ofmt\@firstofone
 % avoid brackets around text for \cite:
 \def\@biblabel#1{}
 \def\@cite#1#2{{#1\if@tempswa , #2\fi}}
\makeatother
\newlength{\cslhangindent}
\setlength{\cslhangindent}{1.5em}
\newlength{\csllabelwidth}
\setlength{\csllabelwidth}{3em}
\newenvironment{CSLReferences}[2] % #1 hanging-indent, #2 entry-spacing
 {\begin{list}{}{%
  \setlength{\itemindent}{0pt}
  \setlength{\leftmargin}{0pt}
  \setlength{\parsep}{0pt}
  % turn on hanging indent if param 1 is 1
  \ifodd #1
   \setlength{\leftmargin}{\cslhangindent}
   \setlength{\itemindent}{-1\cslhangindent}
  \fi
  % set entry spacing
  \setlength{\itemsep}{#2\baselineskip}}}
 {\end{list}}
\usepackage{calc}
\newcommand{\CSLBlock}[1]{\hfill\break\parbox[t]{\linewidth}{\strut\ignorespaces#1\strut}}
\newcommand{\CSLLeftMargin}[1]{\parbox[t]{\csllabelwidth}{\strut#1\strut}}
\newcommand{\CSLRightInline}[1]{\parbox[t]{\linewidth - \csllabelwidth}{\strut#1\strut}}
\newcommand{\CSLIndent}[1]{\hspace{\cslhangindent}#1}
%\usepackage[svgnames]{xcolor}
%\definecolor{codebackground}{RGB}{240, 240, 235}
%\definecolor{codebackground}{RGB}{117, 128, 124}
%\AtBeginDocument{\colorlet{defaultcolor}{.}}
%\definecolor{bg}{HTML}{282828} % from https://github.com/kevinsawicki/monokai
%\usepackage[outputdir=build]{minted}
%\setminted{style=monokai,bgcolor=bg}
%\setmintedinline{style=monokai,bgcolor=None}
%\definecolor{Text}{HTML}{F8F8F2}
%\AddToHook{cmd/mintinline/before}{\color{Text}}
%\AddToHook{cmd/mintinline/after}{}
    %\AtBeginEnvironment{minted}{\color{Text}}
\usepackage{pgfornament}
\usepackage{setspace}
\usepackage{microtype}
\usepackage{fontspec}
    \defaultfontfeatures{Numbers=OldStyle}
    \setmainfont{STIX Two Text}
\setmonofont{PragmataPro Mono Liga}
%\renewcommand{\footnote}[1]{\sidenote{#1}}
%\renewcommand{\familydefault}{\sfdefault}
\fancyfoot[LEO]{\footnotesize © 2026 Andrew Hayes.This work (apart from any source code it contains) is licensed under CC BY 4.0. To view a copy of this license, visit https://creativecommons.org/licenses/by/4.0/.}
\ifLuaTeX
  \usepackage{selnolig}  % disable illegal ligatures
\fi
\usepackage{bookmark}
\IfFileExists{xurl.sty}{\usepackage{xurl}}{} % add URL line breaks if available
\urlstyle{same}
\hypersetup{
  pdftitle={Adding Semantic Structure to Unstructured PDFs via the Unix Command Line},
  pdfauthor={Andrew J. Hayes},
  colorlinks=true,
  linkcolor={teal},
  filecolor={Maroon},
  citecolor={Blue},
  urlcolor={blue},
  pdfcreator={LaTeX via pandoc}}

\title{Adding Semantic Structure to Unstructured PDFs via the Unix
Command Line}
\author{Andrew J. Hayes}
\date{May 16, 2026}

\begin{document}
\maketitle

\ifdefined\soulregister

\soulregister\MakeTextUppercase{1} \soulregister\MakeTextLowercase{1}
\soulregister\newlinetospace{1} \fi

\doublespacing

\section{Introduction and Use Case}\label{introduction-and-use-case}

This paper demonstrates the concept and simple implementation for a
workflow that an individual scholar may use to enhance the utility of a
private research corpus of ancient or medieval sources in pdf format by
adding a semantically structured document outline to them. Traditional
papers and books are usually divided into sections, each with a header,
and (in the case of books) often furnished with some kind of table of
contents to facilitate rapid navigation to a particular point of
interest in that structure. The pdf standard, which intends to codify a
digital representation of paper documents that would be portable across
different operating systems,\footnote{\emph{ISO 32000-2 Portable
  Document Format -- Part 2: PDF 2.0}, 2020, viii (§0.1); Jonas
  Gamalielsson and Björn Lundell, \emph{Experiences from Implementing
  PDF in Open Source: Challenges and Opportunities for Standardisation
  Processes}, vol. 2013 8th International Conference on Standardization
  and Innovation in Information Technology (SIIT) (IEEE, 2013).} has
long supported a digital implementation of this feature: a navigable
table of contents, usually in a sidebar. But, despite the fact that
papers, books, and theses are now commonly accessed in pdf form, they
often lack such an outline. This paper shows how to automate the process
of adding one using Unix command line tools.

This workflow will only be useful if the document in question is
composed of searchable digital text, either because it was born digital,
or because optical character recognition (OCR) has subsequently added a
text layer to each page. It is not suited for for very short texts or
texts that inherently lack a discernible structure or possess it only
minimally, such as scans of medieval manuscripts. The most usual
application will be to critical editions of primary sources or to the
translation volumes that accompany them. A logical secondary application
would be to longer articles with many sections.

This presentation aims to show what is possible using open source Python
tooling and to provide a simple implementation of command line scripts
that the reader may use and adapt for his or her own specific needs.
Along the same lines, the paper includes a few observations about what
was needed to implement the scripts in Python and POSIX shell, with the
hope that this will help others learn enough to adapt the scripts to
their purposes and avoid pitfalls. This paper and the accompanying
source code will be released under an open source license. Code may be
updated in the future to account for changes in the Python ecosystem, or
the Docling library on which the code depends. The workflow demonstrated
here is compatible with the one demonstrated in the author's paper:
``Efficient Vocabulary Discovery in Late Antique Texts'' given at the
2025 Annual Meeting of the North American Patristics Society, and
\href{https://github.com/theusefulmonk/Vocabulary-Discovery-in-Late-Antique-Texts}{publicly
available here}.\footnote{\textbf{Note on AI Use:} Generative AI was not
  used to research the ideas and techniques contained in, nor produce
  the text of, this paper or its associated source code.}

\section{Brief Conceptual Background to the
Tooling}\label{brief-conceptual-background-to-the-tooling}

Python is a general purpose interpreted programming language widely used
in scientific research and in academia more generally. It enjoys a rich
standard library, and can be used (among many other applications) to
write composable command line scripts, following the Unix software tools
approach.\footnote{Brian W. Kernighan and P. J. Plauger, \emph{Software
  Tools} (Reading, Massachusetts: Addison-Wesley Professional, 1976),
  1--6.} According to this approach tools should ideally perform a
single task well, but be designed to receive and pass along input and
output in a standard way, to enable composition. One achieves such
composition by interleaving calls to specific commands or scripts with
the pipe character (\texttt{\textbar{}}). The output of one program
becomes the input of the next in the pipeline until the final desired
result is output.\footnote{Richard Blum and Christine Bresnahan,
  \emph{Linux Command Line and Shell Scripting Bible} (John Wiley \&
  Sons, 2015), 279--84.} The workflow demonstrated in this paper uses
the Python standard library in conjunction with Docling to script a
pipeline that takes a pdf lacking a digital Document outline, creates
such an outline by parsing the document's section headers, adjusts it,
and applies it to a new copy of the pdf.

Docling is relatively recent open source python library and command line
tool, developed by IBM for parsing pdfs and other human-readable
documents\footnote{{``Docling''}
  (\url{https://docling-project.github.io/docling/}, n.d.).} so that
their data can be more easily ingested by another application, such as a
large language model and thus augment it with domain-specific knowledge.
This typical use-case is known as RAG (Retrieval Augmented
Generation).\footnote{{``Retrieval-Augmented Generation (RAG) Vs.
  Fine-Tuning,''} The official Red Hat blog
  (\url{https://www.redhat.com/en/topics/ai/rag-vs-fine-tuning}, 2026).}
This is a broad domain, and the Docling library is powerful and
feature-rich. Nevertheless, this workflow uses only a small subset of
its abilities for a comparatively simple task: extracting the names and
locations of section headings in a structured form that can be
manipulated programmatically.

Pdfcpu is an open source library and command line tool written in Go
which is able to perform a broad range of pdf manipulations.\footnote{{``Pdfcpu:
  A PDF Processor Written in Go''} (\url{https://pdfcpu.io}, n.d.).} In
this workflow it is used to apply the tree structure of the document
outline to the pdf, as the final stage in the workflow.

It is important to note that the terminology for the pdf document
outline varies. The current PDF ISO standard describes it this way:

\begin{quote}
The outline consists of a tree-structured hierarchy of outline items
(sometimes called bookmarks), which serve as a visual table of contents
to display the document's structure to the user.\footnote{\emph{ISO
  32000-2 Portable Document Format -- Part 2}, sec. 12.3.3.}
\end{quote}

The wording of the standard suggests that the canonical name for the
items in the hierarchy is ``outline items.'' Adobe Acrobat's user
interface at the time of writing refers to the items in the structure as
bookmarks. Pdfcpu's documentation also employs this terminology. Note
that bookmarks need not form part of a tree. They can take the form of
navigation targets in a flat list. In this paper such bookmarks or
document outlines and their outline items will also be referred to
generically as tables of contents.

\section{Workflow}\label{workflow}

\textbf{Cautionary Note:}

When working with any important file, such as laboriously obtained scans
of physical originals, it is good practice always to edit copies rather
than overwriting the original files.

The workflow consists of three phases: generating, updating, and
applying the table of contents. Below, I offer an outline of these three
phases, followed by some practical considerations for real-world use.

\subsection{Phase One: Generate the Table of
Contents}\label{phase-one-generate-the-table-of-contents}

In order to add a table of contents to a pdf, one first has to derive
the information from the file and represent it in a structured way. To
do this, we use Docling's ability to process a pdf and export a stream
of JSON\footnote{{``Introducing JSON''}
  (\url{https://www.json.org/json-en.html}, n.d.); \emph{Standard ECMA
  404: The JSON Data Interchange Syntax}, 2017.} (which can also, if
desired, be saved to a file). This JSON contains Docling's
``understanding'' of the whole document, but to create a document
outline we need to extract only the information about section headings
from the JSON Docling produces.

At least two additional practical considerations are relevant. First,
calls to Docling are time consuming and computationally expensive. For
this reason, the script caches the result of a previous run. A full
re-parse of the document can be forced by invoking the \texttt{tcgen}
script with the \texttt{-B} flag. Second, since medievalists and
humanists more generally often have to work with texts in languages such
as Latin or Greek, it is important to output the json text in utf-8
encoding. In the script, this is done by ensuring that calls to Python's
\texttt{json.dumps} (dump string) function are passed the
\texttt{ensure\_ascii=False} parameter. Users of the script must, of
course, have a unicode-capable terminal emulator and the appropriate
fonts installed.

In our example, we will use a born-digital pdf critical edition of the
letters of the \emph{Corpus Dionysiacum}.\footnote{Pseudo-Dionysius
  Areopagita, \emph{Corpus Dionysiacum, II}, ed. Günther Heil and Adolf
  Martin Ritter, vol. 36 (Berlin: Walther de Gruyter, 1991).} Note that,
due to copyright, it cannot be distributed in the repository along with
this paper. Readers are encouraged to supply their own pdf file(s) for
experimentation. Despite being purchased as a natively digital pdf, it
lacks any document outline. It contains the usual sorts of headers that
a printed critical edition might contain: an introductory section for
\emph{sigla} used in the edition, followed by the critical text of each
letter, whose titles, like the letters themselves, are printed in Greek.
To generate a table of contents, we simply invoke the script with the
filename as an argument:

\begin{Shaded}
\begin{Highlighting}[]
\ExtensionTok{./tcgen}\NormalTok{ ./sources/Pseudo{-}Dionysius\_Areopagita\_1991\_416.pdf}
\end{Highlighting}
\end{Shaded}

This outputs a json string containing an array, each element of which is
a json object containing a title key and a page key, indicating the
title of the section heading and the absolute page number of the pdf on
which it occurs.

The resulting json representation of the table of contents is not easy
to read because it is visually compressed. As a result, it is hard to
check for any errors in the result. This can be fixed by piping the
result into \texttt{./tcedit} without any flags:

\begin{Shaded}
\begin{Highlighting}[]
\ExtensionTok{./tcgen}\NormalTok{ ./sources/Pseudo{-}Dionysius\_Areopagita\_1991\_416.pdf }\KeywordTok{|} \ExtensionTok{./tcedit}
\end{Highlighting}
\end{Shaded}

Visually inspecting this output leads us to the next phase of the
workflow.

\subsection{Phase Two: Update the Table of
Contents}\label{phase-two-update-the-table-of-contents}

\begin{Shaded}
\begin{Highlighting}[]
\ErrorTok{1.} \FunctionTok{\{}\ErrorTok{\textquotesingle{}title\textquotesingle{}}\FunctionTok{:} \ErrorTok{\textquotesingle{}EPISTULAE\textquotesingle{}}\FunctionTok{,} \ErrorTok{\textquotesingle{}page\textquotesingle{}}\FunctionTok{:} \DecValTok{1}\FunctionTok{\}}
\ErrorTok{2.} \FunctionTok{\{}\ErrorTok{\textquotesingle{}title\textquotesingle{}}\FunctionTok{:} \ErrorTok{\textquotesingle{}SIGLA}  \ErrorTok{CODICUM\textquotesingle{}}\FunctionTok{,} \ErrorTok{\textquotesingle{}page\textquotesingle{}}\FunctionTok{:} \DecValTok{2}\FunctionTok{\}}
\ErrorTok{3.} \FunctionTok{\{}\ErrorTok{\textquotesingle{}title\textquotesingle{}}\FunctionTok{:} \ErrorTok{\textquotesingle{}ΕΠΙΣΤΟΛΑΙ} \ErrorTok{ΔΙΑΦΟΡΟΙ\textquotesingle{}}\FunctionTok{,} \ErrorTok{\textquotesingle{}page\textquotesingle{}}\FunctionTok{:} \DecValTok{5}\FunctionTok{\}}
\ErrorTok{4.} \FunctionTok{\{}\ErrorTok{\textquotesingle{}title\textquotesingle{}}\FunctionTok{:} \ErrorTok{\textquotesingle{}α} \ErrorTok{ΓΑΙΩΙ}  \ErrorTok{ΘΕΡΑΠΕΥΤΗI\textquotesingle{}}\FunctionTok{,} \ErrorTok{\textquotesingle{}page\textquotesingle{}}\FunctionTok{:} \DecValTok{6}\FunctionTok{\}}
\ErrorTok{5.} \FunctionTok{\{}\ErrorTok{\textquotesingle{}title\textquotesingle{}}\FunctionTok{:} \ErrorTok{\textquotesingle{}ß} \ErrorTok{ΤΩΙ}  \ErrorTok{ΑΥΤΩΙ\textquotesingle{}}\FunctionTok{,} \ErrorTok{\textquotesingle{}page\textquotesingle{}}\FunctionTok{:} \DecValTok{8}\FunctionTok{\}}
\ErrorTok{6.} \FunctionTok{\{}\ErrorTok{\textquotesingle{}title\textquotesingle{}}\FunctionTok{:} \ErrorTok{\textquotesingle{}y\textquotesingle{}}\FunctionTok{,} \ErrorTok{\textquotesingle{}page\textquotesingle{}}\FunctionTok{:} \DecValTok{9}\FunctionTok{\}}
\ErrorTok{7.} \FunctionTok{\{}\ErrorTok{\textquotesingle{}title\textquotesingle{}}\FunctionTok{:} \ErrorTok{\textquotesingle{}}\DecValTok{5}\ErrorTok{\textquotesingle{}}\FunctionTok{,} \ErrorTok{\textquotesingle{}page\textquotesingle{}}\FunctionTok{:} \DecValTok{10}\FunctionTok{\}}
\ErrorTok{8.} \FunctionTok{\{}\ErrorTok{\textquotesingle{}title\textquotesingle{}}\FunctionTok{:} \ErrorTok{\textquotesingle{}ΤΩΙ}  \ErrorTok{ΑΥΤΩΙ\textquotesingle{}}\FunctionTok{,} \ErrorTok{\textquotesingle{}page\textquotesingle{}}\FunctionTok{:} \DecValTok{10}\FunctionTok{\}}
\ErrorTok{9.} \FunctionTok{\{}\ErrorTok{\textquotesingle{}title\textquotesingle{}}\FunctionTok{:} \ErrorTok{\textquotesingle{}ε} \ErrorTok{ΔΩΡΟΘΕΩΙ}  \ErrorTok{ΛΕΙΤΟΥΡΓΩΙ\textquotesingle{}}\FunctionTok{,} \ErrorTok{\textquotesingle{}page\textquotesingle{}}\FunctionTok{:} \DecValTok{12}\FunctionTok{\}}
\ErrorTok{10.} \FunctionTok{\{}\ErrorTok{\textquotesingle{}title\textquotesingle{}}\FunctionTok{:} \ErrorTok{\textquotesingle{}S} \ErrorTok{ΣΩΠΑΤΡΩΙ}  \ErrorTok{IEPEI\textquotesingle{}}\FunctionTok{,} \ErrorTok{\textquotesingle{}page\textquotesingle{}}\FunctionTok{:} \DecValTok{14}\FunctionTok{\}}
\ErrorTok{11.} \FunctionTok{\{}\ErrorTok{\textquotesingle{}title\textquotesingle{}}\FunctionTok{:} \ErrorTok{\textquotesingle{}ι} \ErrorTok{ΠΟΛΥΚΑΡΠΩΙ}  \ErrorTok{ΙΕΡΑΡΧΗΙ\textquotesingle{}}\FunctionTok{,} \ErrorTok{\textquotesingle{}page\textquotesingle{}}\FunctionTok{:} \DecValTok{15}\FunctionTok{\}}
\ErrorTok{12.} \FunctionTok{\{}\ErrorTok{\textquotesingle{}title\textquotesingle{}}\FunctionTok{:} \ErrorTok{\textquotesingle{}η} \ErrorTok{ΔΗΜΟΦΙΛΩΙ}  \ErrorTok{ΘΕΡΑΠΕΥΤΗ} \ErrorTok{I\textquotesingle{}}\FunctionTok{,} \ErrorTok{\textquotesingle{}page\textquotesingle{}}\FunctionTok{:} \DecValTok{21}\FunctionTok{\}}
\ErrorTok{13.} \FunctionTok{\{}\ErrorTok{\textquotesingle{}title\textquotesingle{}}\FunctionTok{:} \ErrorTok{\textquotesingle{}θ\textquotesingle{}}\FunctionTok{,} \ErrorTok{\textquotesingle{}page\textquotesingle{}}\FunctionTok{:} \DecValTok{43}\FunctionTok{\}}
\ErrorTok{14.} \FunctionTok{\{}\ErrorTok{\textquotesingle{}title\textquotesingle{}}\FunctionTok{:} \ErrorTok{\textquotesingle{}ΤΙΤΩΙ}  \ErrorTok{ΊΕΡΑΡΧΗΙ\textquotesingle{}}\FunctionTok{,} \ErrorTok{\textquotesingle{}page\textquotesingle{}}\FunctionTok{:} \DecValTok{43}\FunctionTok{\}}
\ErrorTok{15.} \FunctionTok{\{}\ErrorTok{\textquotesingle{}title\textquotesingle{}}\FunctionTok{:} \ErrorTok{\textquotesingle{}ΊΩΑΝΝΗΙ}  \ErrorTok{ΘΕΟΛΟΓΩΙ}\FunctionTok{,}  \ErrorTok{ΆΠΟΣΤΟΛΩΙ}  \ErrorTok{ΚΑΙ}  \ErrorTok{ΕΥΑΓΓΕΛΙΣΤΗΙ} \ErrorTok{ΠΕΡΙΟΡΙΣΘΕΝΤΙ}  \ErrorTok{ΚΑΤΑ} \ErrorTok{ΠΑΤΜΟΝ}  \ErrorTok{ΤΗΝ}  \ErrorTok{ΝΗΣΟΝ\textquotesingle{}}\FunctionTok{,} \ErrorTok{\textquotesingle{}page\textquotesingle{}}\FunctionTok{:} \DecValTok{58}\FunctionTok{\}}
\end{Highlighting}
\end{Shaded}

The output of \texttt{tcgen} is imperfect. It contains some
(semantically) duplicate nodes, because Docling treated some headers as
two separate headers on the same page. It also contains some basic
labelling mistakes. How can we more easily spot them? \texttt{tcedit}
offers a way.

The script \texttt{tcedit} receives the json stream from the previous
command (here \texttt{tcgen}). When given no options, it outputs an
enumerated list of bookmark nodes so that the user can inspect them and
decide what changes might be necessary. This intermediate representation
is intended for the user and not for passing on to another command in
the pipeline. In this case, if we invoke it with the same filename as
before:

\begin{Shaded}
\begin{Highlighting}[]
\ExtensionTok{./tcedit}\NormalTok{ ./sources/Pseudo{-}Dionysius\_Areopagita\_1991\_416.pdf}
\end{Highlighting}
\end{Shaded}

we will see the enumerated list of bookmark nodes. This reveals the
duplicated nodes, as well as nodes that mistakenly treated Greek text
like Roman script.

It can be invoked with the \texttt{-d} (or in long form
\texttt{-\/-delete}) flag, along with the index numbers of the nodes to
delete. For example, to delete the seventh and the thirteenth bookmark,
one could invoke it as follows:

\begin{Shaded}
\begin{Highlighting}[]
\ExtensionTok{./tcgen}\NormalTok{ ./sources/Pseudo{-}Dionysius\_Areopagita\_1991\_416.pdf }\KeywordTok{|} \DataTypeTok{\textbackslash{}}
\ExtensionTok{./tcedit} \AttributeTok{{-}d}\NormalTok{ 7 13}
\end{Highlighting}
\end{Shaded}

It will output the json string to standard output with the requested
nodes deleted.

Using the \texttt{-u} (long form \texttt{-\/-update}) flag, it is
possible to make changes to the text of a given node using a KEY=VALUE
syntax in which the VALUE is itself a tuple consisting of the portion of
the node that needs editing. For example:

\begin{Shaded}
\begin{Highlighting}[]
\ExtensionTok{./tcgen}\NormalTok{ sources/Pseudo{-}Dionysius\_Areopagita\_1991\_416.pdf }\KeywordTok{|} \DataTypeTok{\textbackslash{}}
\ExtensionTok{./tcedit} \AttributeTok{{-}u}\NormalTok{ 6=}\StringTok{"title,γ ΤΩΙ ΑΥΤΩΙ"} \DataTypeTok{\textbackslash{}}
\NormalTok{8=}\StringTok{"title,δ ΤΩΙ ΑΥΤΩΙ"} \DataTypeTok{\textbackslash{}}
\NormalTok{14=}\StringTok{"title,θ ΤΙΤΩΙ ᾽ΙΕΡΑΡΧΗΙ"} \DataTypeTok{\textbackslash{}}
\NormalTok{{-}d 7 13}
\end{Highlighting}
\end{Shaded}

This pipeline updates the title portions of nodes 6, 8, and 14. Both the
\texttt{-u} and \texttt{-d} flags may be used together with the same
command. The order in which they are given does not matter.
\texttt{tcedit} will always perform update operations before the delete
operations in a single invocation. Once we are satisfied that the edits
are correct, the \texttt{-f} (long form \texttt{-\/-final}) flag can be
added to output it in the form needed for Pdfcpu to apply it to the
finished pdf.

\subsection{Phase Three: Apply the Table of Contents to the PDF
file}\label{phase-three-apply-the-table-of-contents-to-the-pdf-file}

The final stage of the workflow is to apply the output of the
\texttt{tcedit} script to a copy of the original pdf file and output it.
A simple shell script, \texttt{tcApply.sh} has been provided for this
purpose:

\begin{Shaded}
\begin{Highlighting}[]
\ExtensionTok{./tcgen}\NormalTok{ sources/Pseudo{-}Dionysius\_Areopagita\_1991\_416.pdf }\KeywordTok{|} \DataTypeTok{\textbackslash{}}
\ExtensionTok{./tcedit} \AttributeTok{{-}u}\NormalTok{ 6=}\StringTok{"title,γ ΤΩΙ ΑΥΤΩΙ"} \DataTypeTok{\textbackslash{}}
\NormalTok{8=}\StringTok{"title,δ ΤΩΙ ΑΥΤΩΙ"} \DataTypeTok{\textbackslash{}}
\NormalTok{14=}\StringTok{"title,θ ΤΙΤΩΙ ᾽ΙΕΡΑΡΧΗΙ"} \DataTypeTok{\textbackslash{}}
\NormalTok{{-}d 7 13 }\AttributeTok{{-}f} \KeywordTok{|} \DataTypeTok{\textbackslash{}}
\ExtensionTok{./tcApply.sh}\NormalTok{ sources/Pseudo{-}Dionysius\_Areopagita\_1991\_416.pdf}
\end{Highlighting}
\end{Shaded}

Note the use of the \texttt{-f} flag to finish outputting the result of
\texttt{tcedit}, which is then piped into the \texttt{tcApply.sh}
script, which produces a file called simply \texttt{output.pdf}. This
file contains the navigable table of contents. It is now ready to use.
Ordinarily, one would then meaningfully rename the file and store it in
a personal digital library of texts.

\subsection{Considerations in Real-World
Use}\label{considerations-in-real-world-use}

On the Unix command line, it is often useful to compose a pipeline with
repeated invocations of the same command but with different flags or
arguments. Doing this makes it easy to experiment and build up
incrementally the desired result. Because the source is never modified
directly, one can always repeat new variations of the pipeline until the
desired result is achieved, which can then be written to a file.
\texttt{tcedit} largely conforms to this usage pattern. It can be
applied multiple times in the pipeline, breaking up a complex invocation
into simple parts. It can always be added to the chain without a flag to
inspect the output at that point in the transformation process.

It may also happen that one needs to stop work in the midst of an
exploratory session in which one has not completed all desired
transformations. In this case, you can redirect the bookmarks object to
a file (equivalent to issuing a ``Save'' command via a GUI). It is
possible to reload this file later to resume the transformation process.
Here is what that would look like:

\begin{Shaded}
\begin{Highlighting}[]
\ExtensionTok{./tcgen}\NormalTok{ sources/Pseudo{-}Dionysius\_Areopagita\_1991\_416.pdf }\KeywordTok{|}\DataTypeTok{\textbackslash{}}
\ExtensionTok{./tcedit} \AttributeTok{{-}u}\NormalTok{ 6=}\StringTok{"title,γ ΤΩΙ ΑΥΤΩΙ"} \DataTypeTok{\textbackslash{}}
\NormalTok{8=}\StringTok{"title,δ ΤΩΙ ΑΥΤΩΙ"} \DataTypeTok{\textbackslash{}}
\NormalTok{14=}\StringTok{"title,θ ΤΙΤΩΙ ᾽ΙΕΡΑΡΧΗΙ"} \DataTypeTok{\textbackslash{}}
\NormalTok{{-}d 7 13 }\OperatorTok{\textgreater{}}\NormalTok{ unfinished.json}
\end{Highlighting}
\end{Shaded}

This would save the bookmark list to the file \texttt{unfinished.json}.
It could easily be picked up again using the standard Unix tool
\texttt{cat}:

\begin{Shaded}
\begin{Highlighting}[]
\FunctionTok{cat}\NormalTok{ unfinished.json }\KeywordTok{|} \ExtensionTok{./tcedit} \AttributeTok{{-}u} \DataTypeTok{\textbackslash{}}
\NormalTok{13=}\StringTok{"title,ι ᾽ΙΩΑΝΝΗΙ ΘΕΟΛΟΓΩΙ, ΆΠΟΣΤΟΛΩΙ ΚΑΙ ΕΥΑΓΓΕΛΙΣΤΗΙ ΠΕΡΙΟΡΙΣΘΕΝΤΙ ΚΑΤΑ ΠΑΤΜΟΝ ΤΗΝ ΝΗΣΟΝ"} \KeywordTok{|} \ExtensionTok{./tcedit}
\end{Highlighting}
\end{Shaded}

Here, the contents of the \texttt{unfinished.json} file are piped back
into the \texttt{tcedit} script to complete one more update. The last
node in the bookmark list does not begin with a lowercase Greek letter.
All the other section headers begin with their appropriate Greek letter
in sequence. But with this final transformation accomplished, the last
section header matches the format of all the rest.

Docling's ability to output pdf structure as json is a significant
advantage because json is so widely used and widely supported. In fact,
for more advanced transformation, it is possible to dispense with the
\texttt{tcedit} script and use \texttt{jq}, a robust command line tool
for filtering json.\footnote{{``/Jq''} (\url{https://jqlang.org}, n.d.).}
Just like any other Unix filter, it can be added to the pipeline to
achieve whatever transformation is desired. \texttt{jq} is out of scope
for this tutorial, but it shows what is possible with composable command
line tools.

A digital humanist could easily modify and extend these scripts. One
could, for example, create a version of these tools that would output
just the text of each heading with or without page numbers. This could
be redirected to another file and turned into the skeleton of a handout
for a talk. This output could in turn be modified or piped into pandoc
to produce the final handout. Docling is not limited to ingesting pdfs.
It can also ingest \LaTeX{} files and MS Word files.

\section{Limitations and Summary}\label{limitations-and-summary}

The workflow offered here makes the most sense as a tool for an
individual scholar or a small team with limited budget. Alternatives
exist. One could use a variety of commercially available GUI software
tools to edit the document outline of an existing pdf. The most obvious
one is Adobe Acrobat, distributed by the company that developed the PDF
standard originally. PDFExpert is a powerful (albeit Mac only)
alternative. Xodo distributes a free pdf reader that can edit document
outlines. Other commercially available software, including web-based
tools may have similar or even more advanced capabilities. But of
course, commercial solutions cost money and cannot be freely adapted
like open source software can. Perhaps more importantly for efficient
processing of large personal libraries, they are more complex to
automate or cannot be automated at all. The primary advantage that
command line tools provide is that Docling can generate a reasonable
approximation whose results can be processed in a batch or in an
automated fashion. Creating a document outline for a large critical
edition by hand using a GUI tool, however powerful, would be laborious
and error prone.

The scripts offered in this tutorial are proof-of-concept. They have
immediate utility, but also a few key limitations that one might want to
improve upon in a version adapted for personal use. Most notably,
\texttt{tcedit} lacks the ability to perform more complex
transformations beyond simple node deletion and editing. It cannot nest
them as sub-headings. This limitation is also reflected in Docling,
which does not seem to be able to reliably interpret the section header
hierarchy. In the interest of speed, the \texttt{tcgen} script does not
perform OCR directly, though this is a capability that Docling offers,
and it could be added by adapting the \texttt{tcgen} script and
installing OCR dependencies on the local system.

Docling allows the creation of powerful command line tools in Python for
extracting the information needed for a pdf document outline. Pdfcpu
makes it easy to apply this document outline to a file. The result is a
more useful digital text. The Unix Software Tools philosophy makes this
utility possible by creating an environment in which commands can be
iteratively composed, experimented with, and automated. And such tools
are available for free with few restrictions on their use.

\clearpage

\section*{Bibliography}\label{bibliography}
\addcontentsline{toc}{section}{Bibliography}

\phantomsection\label{refs}
\begin{CSLReferences}{1}{0}
\bibitem[\citeproctext]{ref-blum2015}
Blum, Richard, and Christine Bresnahan. \emph{Linux Command Line and
Shell Scripting Bible}. John Wiley \& Sons, 2015.

\bibitem[\citeproctext]{ref-docling2026}
{``Docling.''} \url{https://docling-project.github.io/docling/}, n.d.

\bibitem[\citeproctext]{ref-gamalielsson2013}
Gamalielsson, Jonas, and Björn Lundell. \emph{Experiences from
Implementing PDF in Open Source: Challenges and Opportunities for
Standardisation Processes}. Vol. 2013 8th International Conference on
Standardization and Innovation in Information Technology (SIIT). IEEE,
2013.

\bibitem[\citeproctext]{ref-json}
{``Introducing JSON.''} \url{https://www.json.org/json-en.html}, n.d.

\bibitem[\citeproctext]{ref-pdfassociation2020}
\emph{ISO 32000-2 Portable Document Format -- Part 2: PDF 2.0}, 2020.

\bibitem[\citeproctext]{ref-jq}
{``./Jq.''} \url{https://jqlang.org}, n.d.

\bibitem[\citeproctext]{ref-kernighan1976}
Kernighan, Brian W., and P. J. Plauger. \emph{Software Tools}. Reading,
Massachusetts: Addison-Wesley Professional, 1976.

\bibitem[\citeproctext]{ref-pdfcpu}
{``Pdfcpu: A PDF Processor Written in Go.''} \url{https://pdfcpu.io},
n.d.

\bibitem[\citeproctext]{ref-dionysius1991}
Pseudo-Dionysius Areopagita. \emph{Corpus Dionysiacum, II}. Edited by
Günther Heil and Adolf Martin Ritter. Vol. 36. Berlin: Walther de
Gruyter, 1991.

\bibitem[\citeproctext]{ref-redhat2026}
The official Red Hat blog. {``Retrieval-Augmented Generation (RAG) Vs.
Fine-Tuning.''}
\url{https://www.redhat.com/en/topics/ai/rag-vs-fine-tuning}, 2026.

\bibitem[\citeproctext]{ref-ecma2017}
\emph{Standard ECMA 404: The JSON Data Interchange Syntax}, 2017.

\end{CSLReferences}

\end{document}
